\documentclass[a4paper]{report}
\usepackage{anysize}
\marginsize{2cm}{2cm}{2cm}{2cm}
\usepackage{pdfpages}
\usepackage[utf8]{inputenc}
\usepackage{xspace}
\usepackage{mathtools}
\usepackage[francais]{babel}
\usepackage{listingsutf8}
\usepackage{listings}
% addresse du tuto pour activer la coloration du code
% http://blog.hikoweb.net/index.php?post/2012/02/04/Colorer-du-code-source-dans-un-rapport-LaTeX-avec-Minted
\usepackage{minted}
%\usepackage{graphics}
\usepackage{graphicx}
\usepackage[framemethod=tikz]{mdframed}
\graphicspath{ {./img/} }
\DeclareUnicodeCharacter{B0}{\textdegree}
\renewcommand{\contentsname}{Sommaire}
%% \usepackage{natbib}
\usepackage{hyperref}
\hypersetup{
    bookmarks=true,         % show bookmarks bar?
    unicode=false,          % non-Latin characters in Acrobat’s bookmarks
    pdfnewwindow=true,      % links in new window
    colorlinks=true,       % false: boxed links; true: colored links
    linkcolor=black,        % color of internal links (change box color with linkbordercolor)
    citecolor=magenta,        % color of links to bibliography
    filecolor=magenta,      % color of file links
    urlcolor=blue           % color of external links
}
\usepackage{breakurl}
\usepackage[T1]{fontenc}
\usepackage[nottoc, notlof, notlot]{tocbibind}
\usepackage{fancyhdr}
\usepackage{lastpage} 

\pagestyle{fancy}


\lhead{}
\chead{}
\rhead{} 

\renewcommand{\headrulewidth}{0pt}

\renewcommand{\footrulewidth}{0.4pt}

\lfoot{System On Chip OpenSources}

\cfoot{}

\rfoot{Page \thepage}

\begin{document}

\begin{titlepage}
  \centering
  \vspace{7\baselineskip}
  {\huge Compte-rendu de projet }\\
  \vspace{1cm}
  {\Huge{\bf Développement d'un System-On-Chip open-source }}\\
  \vspace{1cm}
  Encadrants :\\
  {\bf Fernand LONE SANG}\\
  {\bf Vincent NICOMETTE}\\
  
  \par
  \normalfont
  \vspace{1cm}
  DUZAN Luc\\
  \vspace{0.3cm}
  FARACHE Gabriel\\
  \vspace{0.3cm}
  LONGO Matthieu\\
  \vspace{0.3cm}
  MICHAUD Clément\\
  \vspace{1cm}
  \par
  {\tt duzan@etud.insa-toulouse.fr}\\
  {\tt farache@etud.insa-toulouse.fr}\\
  {\tt longo@etud.insa-toulouse.fr}\\
  {\tt cmichaud@etud.insa-toulouse.fr}\\
  \vspace{1cm}
  4IR 2012-2013
  \begin{figure}[h]
    \centering
    \includegraphics[scale=0.4]{INSA-logo.eps}
  \end{figure}
  
  \vspace{3cm}
  \vfill {\bf R\'esum\'e du projet}\medskip \\
  \fbox{
    \begin{minipage}{0.9\textwidth}
      Dans notre projet tutoré, nous allons nous intéresser plus particulièrement au
      \textit{Milkymist} SoC, qui est un SoC \textit{open-source}, c'est à dire que le code HDL qui le décrit
      est libre de droit.  Ce SoC a été conçu pour être programmé sur un FPGA spécifique le
      Virtex-4 XC4VLX25.  Notre but est d'adapter le code de ce SoC pour le rendre
      compatible avec le Xilinx Spartan6 qui est utilisé à l'INSA de Toulouse.  Nous
      allons aussi ajouter un module à \textit{Milkymist} permettant d'exploiter la propriété de
      certains FPGA qui est de pouvoir être reprogrammé à la volée. C'est-à-dire que nous
      allons permettre au système \textit{Milkymist} de pouvoir ajouter à chaud (sans redémarrer le
      système) des périphériques (à partir de leurs descriptions structurelles) sur le FPGA
      et de les connecter au SoC.
    \end{minipage} 
  }
\end{titlepage}

\tableofcontents

\chapter*{Introduction} \addcontentsline{toc}{chapter}{Introduction}
Aujourd'hui, les appareils électroniques que nous connaissons ont tous une ou plusieurs fonctionnalités utiles dans telle ou telle situation. Les appareils de téléphonie, de  réseaux, de télécoms, de traitement d'images, de traitement du son possèdent tous des cartes différentes, chacune d'entre elles étant optimisée pour sa fonction. Nombre d'entre eux ne se limitent d'ailleurs pas à une seule carte mais en nécessite plusieurs.\\

Imaginons maintenant qu'un appareil puisse s'adapter à un besoin à court terme. Imaginez un système qui n'a pas de vocation particulière ou plutôt qui en a une infinité et que vous pouvez choisir celle que vous voulez à un instant donné. Voilà pourquoi il devient intéressant d'étudier les systèmes adaptatifs qui vont s'imposer dans notre vie de tous les jours avec le développement et la  diffusion des puces électroniques reconfigurables et reconfigurables à chaud appelées \fpgas{} pour \english{Field-Programmable Gate Array}.\\

Ce projet est une tentative d'exploitation de cette nouvelle technologie portant le nom de \english{partial reconfiguring} qui permet à une puce \fpga{} de se reconfigurer à chaud, c'est-à-dire pendant qu'une partie du \fpga{} est encore active dans une zone particulière de la puce. Nous devions, après avoir étudier le sujet des FPGA, des \english{Systems-On-Chip} et du \english{partial reconfiguring},  adapter le \english{System-On-Chip} open source \textit{Milkymist}, à la carte \nexys{}.  Le but était ensuite d’y intégrer un module de \english{partial reconfiguring}. Cependant certaines difficultés nous ont amenés à concevoir notre propre \english{System-On-Chip} pour la \nexys{} et à étudier la possibilité de faire communiquer ce système avec une autre carte équipé d’un FPGA Virtex5 sur lequel serait mis en place le module de \english{partial reconfiguring}.\\

Ce rapport a pour but d’expliquer le projet et les technologies utilisées. Pour cela, nous relatons notre stratégie d’organisation dans une première partie. Nous détaillons dans une seconde partie le fonctionnement des \fpgas{} et des \english{Systems-On-Chip}. La troisième partie traite de notre tentative de portage du système \textit{Milkymist} et des difficultés rencontrées. La quatrième partie relate la conception d'un nouveau \english{System-On-Chip} développé par l'équipe. Enfin, La cinquième partie fait un bilan du projet et des connaissances acquises.
%Nous verrons dans un première partie comment l'équipe a été organisé, puis nous présenterons le contexte du projet. Premièrement, nous parlerons de la tentative de portage de MilkyMist et ensuite du nouveau système créé à partir de zéro. Enfin, nous ferons le bilan de l'expérience que nous avons acquise.



\newpage

\chapter{Organisation}
\newpage

\chapter{Les aspects essentiels du projet}
\newpage

\section{Langage de description de circuits logiques}

\subsection{Présentation de Verilog et VHDL}
\input{langage-1} \newpage

\subsection{Les autres HDL existants}
\input{langage-2} \newpage

\subsection{Verilog VS VHDL}
\input{langage-3} \newpage

\subsection{Le HDL de Milkymist : Migen}
\input{langage-4} \newpage

\section{Technologie de fabrication}

\subsection{ASIC}
\input{fpga-asic-1} \newpage

\subsection{FPGA}
\input{fpga-asic-2} \newpage

\section{Le monde de l'\textit{open source} et son avancement sur les FPGA}
\input{opensource-1} \newpage

\chapter{Tentative de portage d'un \textit{System-on-Chip} libre : \textit{Milkymist}}
\newpage

\section{Portage du système}

\subsection{La simulation}
\input{portage-1} \newpage

\subsection{La dure réalité}
\input{portage-2} \newpage

\section{Hardware à tester}

\subsection{La PSRAM}
\input{test-hardware-1} \newpage

\subsection{Le totalité du système}
\input{test-hardware-2} \newpage

\section{La \textit{Partial Reconfiguration} (PR)}

\subsection{Les différents types de PR possibles}
\input{partial-reconfiguration-1} \newpage

\subsection{\textit{PlanAhead} et l'incompatibilité avec la Spartan6}
\input{partial-reconfiguration-2} \newpage

\subsection{La tentative de diff based}
\input{partial-reconfiguration-3} \newpage

\section{Bilan technique}

\subsection{Les raisons de l'échec}
\input{bilan-milkymist-1} \newpage

\subsection{Réorientation vers une nouvelle solution from scratch}
\input{bilan-milkymist-2} \newpage

\chapter{Vers un nouveau SoC}

\section{De nouveaux objectifs}

\subsection{Déportation du PR vers la Virtex5}
\input{nouveaux-objectifs-1} \newpage

\subsection{Linkage de la Virtex5 à la Nexys3}
\input{nouveaux-objectifs-2} \newpage

\section{Réalisations hardware}

\subsection{Rassemblement d'IPs}
\input{realisations-hardware-1} \newpage

\subsection{Création de nouveaux modules}
\input{realisations-hardware-2} \newpage

\section{Développement un noyau ultra léger}

\subsection{Les spécifications}
\input{realisations-software-1} \newpage

\subsection{Le BIOS}
\input{realisations-software-2} \newpage

\subsection{Le kernel}
\input{realisations-software-3} \newpage

\subsection{Une API système simple}
\input{realisations-software-4} \newpage

\subsection{Un pseudo-terminal}
\input{realisations-software-5} \newpage

\subsection{Le driver pour le VGA}
\input{realisations-software-6} \newpage

\subsection{Le driver pour l'Ethernet}
\input{realisations-software-7} \newpage

\chapter*{Bilan final}
\addcontentsline{toc}{chapter}{Bilan final}

\section{Lessons learned}

\subsection{En terme technique}
\input{lesson-learned-1}

\subsection{En terme d'échec}
\input{lesson-learned-2}

\subsection{En terme d'organisation/d'implication}
\input{lesson-learned-3}

\section{Future works}

\subsection{ICAP}
\input{future-works-1}

\subsection{L'intégration du driver Ethernet au kernel}
\input{future-works-2}

\subsection{Établir un protocole de communication entre les cartes}
\input{future-works-3}

\subsection{Tests de crypto soft VS hard}
\input{future-works-4}

%\appendix

% \chapter{Code des fichiers modifiés}

% \section{global.h}
% \inputminted[linenos=true,frame=single,fontsize=\scriptsize]{c}{../projet_destijl/superviseur/global.h}
% \newpage

% \section{global.c}
% \inputminted[linenos=true,frame=single,fontsize=\scriptsize]{c}{../projet_destijl/superviseur/global.c}
% \newpage

% \section{fonctions.h}
% \inputminted[linenos=true,frame=single,fontsize=\scriptsize]{c}{../projet_destijl/superviseur/fonctions.h}
% \newpage

% \section{fonctions.c}
% \inputminted[linenos=true,frame=single,fontsize=\scriptsize]{c}{../projet_destijl/superviseur/fonctions.c}
% \newpage

%  \section{main.c}
% \inputminted[linenos=true,frame=single,fontsize=\scriptsize]{c}{../projet_destijl/superviseur/main.c}

\end{document}

